%!TeX root = main.tex
%!encoding = UTF-8 Unicode

\section{Requirements for the Integrated EMC Measurement Environment}

The measurement platform must fulfill several requirements to enable reproducible and EMC-relevant characterization under realistic operating conditions.
First, it must provide a classical ground reference that represents the inverter chassis and allows physically meaningful coupling paths for high-frequency currents.
Second, the switching behavior of the test setup must be comparable to that of real applications, so that the measured emission spectra remain representative of practical converter operation.
Third, the platform must support the different measurement methodologies introduced in this work, in particular the operation of pulse chains with multiple consecutive switching events at a defined operating point.
Fourth, the setup must include a conducted-EMC reference measurement based on a LISN to allow direct comparison between intrinsic device emissions and standardized system-level quantities.
Fifth, additional noise sources introduced by the test environment must be minimized.
This includes filtering of the supply paths via the LISN, optical transmission of control signals, battery-powered gate-driver supplies, and a linear load implementation based on an air-core inductor.
Finally, the platform must allow stable operation during pulse-chain measurements without requiring active thermal management; therefore, passive cooling has to be sufficient for the intended operating regime.

The design choices discussed in the following section are derived directly from these requirements and are presented in the same logical order.
